\chapter{Gramatyka w pigułce}
\markboth{Gramatyka w pigułce}{Gramatyka w pigułce}
Tabele referencyjne do szybkiego sprawdzenia. Szczegóły i ćwiczenia są w jednostkach; numer w nawiasie to tydzień.

\section*{Rodzajniki i przypadki}
\begin{center}\small
\begin{tabular}{@{}llllll@{}}
\toprule
 & \textsf{\textbf{pytanie}} & \textsf{\textbf{m}} & \textsf{\textbf{f}} & \textsf{\textbf{n}} & \textsf{\textbf{mnoga}}\\
\midrule
Nominativ & kto? co? & \artm{der} / \de{ein} & \artf{die} / \de{eine} & \artn{das} / \de{ein} & \artf{die}\\
Akkusativ & kogo? co? & \artm{den} / \de{einen} & \artf{die} / \de{eine} & \artn{das} / \de{ein} & \artf{die}\\
Dativ & komu? czemu? & \artm{dem} / \de{einem} & \artf{der} / \de{einer} & \artn{dem} / \de{einem} & \artf{den}\,+\,\de{-n}\\
Genitiv & czyj? & \artm{des} / \de{eines}\,+\,\de{-s} & \artf{der} / \de{einer} & \artn{des} / \de{eines}\,+\,\de{-s} & \artf{der}\\
\bottomrule
\end{tabular}
\end{center}

\section*{Zaimki osobowe}
\begin{center}\small
\begin{tabular}{@{}llllllll@{}}
\toprule
Nom & \de{ich} & \de{du} & \de{er} & \de{sie} & \de{es} & \de{wir} & \de{ihr}\\
Akk & \de{mich} & \de{dich} & \de{ihn} & \de{sie} & \de{es} & \de{uns} & \de{euch}\\
Dat & \de{mir} & \de{dir} & \de{ihm} & \de{ihr} & \de{ihm} & \de{uns} & \de{euch}\\
\bottomrule
\end{tabular}
\end{center}
Liczba mnoga \de{sie}: Akk \de{sie}, Dat \de{ihnen}. Grzecznościowe \de{Sie}: Akk \de{Sie}, Dat \de{Ihnen}. (8, 14)

\section*{Czasowniki: odmiana}
\begin{center}\small
\begin{tabular}{@{}lllllll@{}}
\toprule
 & \de{ich} & \de{du} & \de{er/sie/es} & \de{wir} & \de{ihr} & \de{sie/Sie}\\
\midrule
regularne & \de{-e} & \de{-st} & \de{-t} & \de{-en} & \de{-t} & \de{-en}\\
\de{sein} & \de{bin} & \de{bist} & \de{ist} & \de{sind} & \de{seid} & \de{sind}\\
\de{haben} & \de{habe} & \de{hast} & \de{hat} & \de{haben} & \de{habt} & \de{haben}\\
\de{werden} & \de{werde} & \de{wirst} & \de{wird} & \de{werden} & \de{werdet} & \de{werden}\\
\bottomrule
\end{tabular}
\end{center}
Modalne (11): \de{kann, muss, will, darf, soll} bez końcówki dla \de{ich} i \de{er/sie/es}.

\section*{Szyk zdania}
\begin{center}\small
\begin{tabular}{@{}p{0.28\linewidth}p{0.64\linewidth}@{}}
\toprule
\textsf{\textbf{Typ}} & \textsf{\textbf{Reguła i przykład}}\\
\midrule
oznajmujące (2) & czasownik na 2.\ miejscu: \de{Heute arbeite ich zu Hause.}\\
pytanie tak/nie & czasownik na 1.\ miejscu: \de{Arbeitest du heute zu Hause?}\\
pytanie otwarte & słowo pytające, czasownik 2.: \de{Wo arbeitest du?}\\
klamra (11, 15) & \de{Ich muss heute früh aufstehen.} \de{Ich habe gestern gearbeitet.}\\
zdanie podrzędne (18, 25) & czasownik odmieniony na końcu: \de{\dots, weil ich krank bin.}\\
rozkaz (12) & \de{Komm!} \de{Kommen Sie!} \de{Steh auf!}\\
\bottomrule
\end{tabular}
\end{center}

\section*{Przyimki i przypadki}
\begin{center}\small
\begin{tabular}{@{}p{0.2\linewidth}p{0.72\linewidth}@{}}
\toprule
\textsf{\textbf{Przypadek}} & \textsf{\textbf{Przyimki}}\\
\midrule
Akkusativ & \de{durch, für, gegen, ohne, um}\\
Dativ & \de{aus, bei, mit, nach, seit, von, zu, gegenüber}\\
Dativ \textit{lub} Akkusativ & \de{in, auf, an, unter, über, neben, vor, hinter, zwischen}: \de{Wo?} $\rightarrow$ Dat, \de{Wohin?} $\rightarrow$ Akk (10)\\
Genitiv & \de{wegen, trotz, statt, während, innerhalb, außerhalb} (27)\\
\bottomrule
\end{tabular}
\end{center}

\section*{Przymiotniki}
Bez końcówki po \de{sein}. Przed rzeczownikiem: po \de{der/die/das} \de{-e} lub \de{-en}; po \de{ein-} \de{-er} (m Nom), \de{-es} (n Nom/Akk); bez rodzajnika końcówka
rodzajnika (21, 27). Stopniowanie: \de{-er}, \de{am -sten}; nieregularne \de{gut/besser/am besten}, \de{viel/mehr/am meisten}, \de{gern/lieber/am liebsten} (19).

\section*{Czasy w skrócie}
\begin{center}\small
\begin{tabular}{@{}p{0.2\linewidth}p{0.4\linewidth}p{0.3\linewidth}@{}}
\toprule
\textsf{\textbf{Czas}} & \textsf{\textbf{Forma}} & \textsf{\textbf{Użycie}}\\
\midrule
\de{Präsens} & odmiana & teraźniejszość, plan\\
\de{Perfekt} & \de{haben/sein} $+$ PII & mowa o przeszłości (15, 16)\\
\de{Präteritum} & \de{war, hatte, konnte}; forma pisana & \de{sein}, \de{haben}, modalne; tekst (17, 30)\\
\de{Plusquamperfekt} & \de{hatte/war} $+$ PII & wcześniejsze niż inne w przeszłości (30)\\
\de{Futur I} & \de{werden} $+$ Inf. & postanowienie, przypuszczenie (23)\\
\de{Konjunktiv II} & \de{wäre, hätte, würde} $+$ Inf. & uprzejmość, nierealne (23, 29)\\
\de{Konjunktiv I} & \de{sei, habe, werde} & mowa zależna (39)\\
\de{Passiv} & \de{werden} $+$ PII; \de{sein} $+$ PII & proces; stan (28, 40)\\
\bottomrule
\end{tabular}
\end{center}

\section*{Najczęstsze czasowniki nieregularne}
\begin{center}\small
\begin{tabular}{@{}llll@{}}
\toprule
\textsf{\textbf{Bezokolicznik}} & \textsf{\textbf{\de{er}, Präsens}} & \textsf{\textbf{Präteritum}} & \textsf{\textbf{Partizip II}}\\
\midrule
\de{sein} & \de{ist} & \de{war} & \de{gewesen}\\
\de{haben} & \de{hat} & \de{hatte} & \de{gehabt}\\
\de{werden} & \de{wird} & \de{wurde} & \de{geworden}\\
\de{gehen} & \de{geht} & \de{ging} & \de{gegangen}\\
\de{kommen} & \de{kommt} & \de{kam} & \de{gekommen}\\
\de{fahren} & \de{fährt} & \de{fuhr} & \de{gefahren}\\
\de{nehmen} & \de{nimmt} & \de{nahm} & \de{genommen}\\
\de{geben} & \de{gibt} & \de{gab} & \de{gegeben}\\
\de{sehen} & \de{sieht} & \de{sah} & \de{gesehen}\\
\de{lesen} & \de{liest} & \de{las} & \de{gelesen}\\
\de{sprechen} & \de{spricht} & \de{sprach} & \de{gesprochen}\\
\de{essen} & \de{isst} & \de{aß} & \de{gegessen}\\
\de{finden} & \de{findet} & \de{fand} & \de{gefunden}\\
\de{schreiben} & \de{schreibt} & \de{schrieb} & \de{geschrieben}\\
\de{bleiben} & \de{bleibt} & \de{blieb} & \de{geblieben}\\
\de{denken} & \de{denkt} & \de{dachte} & \de{gedacht}\\
\bottomrule
\end{tabular}
\end{center}
